\subsection{Der stärkste neue Selektor verbindet Orientierung und gesamten Randinhalt}\label{sec:chiralselector}
Für dieselbe geeignete Bulk-/Randquelle sei der chirale Zusammenhang $J=\pi(c_L-c_R)/3$ begründet. Für die Vakuum-Intervallentropie ihrer Rand-CFT auf der Linie gilt bei festem Regulator- und Lorentzrahmen~\cite{castroentropy}
\begin{equation}
s=\frac{dS_{\rm Rand}}{d\log\ell}=\frac{c_L+c_R}{6}.
\end{equation}
Dann ergibt sich
\begin{equation}
c_L=3s+\frac{3J}{2\pi},\qquad
c_R=3s-\frac{3J}{2\pi},\qquad
\boxed{s\ge\frac{|J|}{2\pi}.}
\end{equation}
Die Schranke folgt aus Unitarität, $c_L,c_R\ge0$. Ihre Sättigung schließt gegenläufige gaplose CFT-Beiträge aus. Insbesondere
\begin{equation}
J=\frac{8\pi}{3},\qquad s=\frac43
\quad\Longrightarrow\quad (c_L,c_R)=(8,0).
\end{equation}
Mit tatsächlich vorhandenen unitären affinen E8-Strömen folgt dann aus $248k/(k+30)\le8$ und $k\in\Z_{>0}$ das bereits bekannte $k=1$.

Dieser gemeinsame Test ist eine nützliche Verschärfung der Herkunftsprüfung. Seine beiden Eingangswerte wurden hier jedoch nicht aus dem ursprünglichen P1/P2-Rohprozess berechnet. Einsetzen der bereits gewählten E8-Level-1-Zielquelle wäre eine Konsistenzprüfung. Außerdem schließt der Test keine beliebigen zusätzlichen gegappten Sektoren und keine sämtlichen anderen Bulk-Realisierungen aus.

Die ebenfalls genannte lokale Eichkonstruktion von Thorngren, Preskill und Fidkowski ist ein passender konstruktiver Literaturanschluss, aber keine Herleitung dreier Raumdimensionen oder des vollständigen TFPT-Standardmodells~\cite{thorngrendisentanglers}. Sie behält konkrete Voraussetzungen an lokale Schaltungen und eine geeignete gegappte Grenzfläche. Massen, Yukawa und Alpha bleiben die bereits vorhandenen gemeinsamen Anforderungen aus Abschnitt~\ref{sec:physicalconstraints}; sie werden durch diese neuen Orientierungstests nicht ersetzt.

\key{\textbf{Was aus dem Anhang übernommen wird:}
Der Rückanschluss der abstrakten Vierzustandsalgebra an die C-Ladungszeichen ist korrekt. Der Drei-Qubit-Zeuge beweist den Informationsverlust durch reine Paaransichten. Die orientierten Replica-Antworten und der gemeinsame $J/s$-Selektor liefern zusätzliche, bedingte Anforderungen an dieselbe Quelle. Die Carry-Vereinfachung wird ausdrücklich nicht als bereits gelöster Raumlift übernommen.}

\section{Die wichtigste neue Verbindung: Die ursprünglichen 24 Felder erzeugen eine Quelle}\label{sec:source}

\subsection{Die direkte Konstruktion}
Aus dem positiven geraden unimodularen $E_8$-Ladungsgitter und einem kompatiblen Kokzyklus entsteht in der erklärten Level-1-Gitterquantisierung
\begin{equation}
\mathcal H_{E_8}=\widehat{\bigoplus_{\lambda\in E_8}}\mathcal F_\lambda,
\qquad \Omega=\ket0,
\qquad L_0=\frac{\norm{\lambda}^2}{2}+\sum_{j=1}^8\sum_{n>0}nN_{j,n}.
\end{equation}
Die grundlegenden geladenen Vertexantworten haben die Form
\begin{equation}
\left\langle\prod_i V_{\alpha_i}(z_i)\right\rangle
=\delta_{\sum_i\alpha_i,0}\,\varepsilon(\alpha_1,\ldots,\alpha_n)
\prod_{i<j}(z_i-z_j)^{\langle\alpha_i,\alpha_j\rangle}.
\end{equation}
Ladungsneutralität, Kokzyklus und Gitterpaarung bestimmen hier gemeinsam die Antwort. Die Unimodularität begründet die Holomorphie der Gitter-VOA; sie muss nicht als zusätzliche freie Auswahl erneut eingeführt werden. Der Übergang einer geeigneten unitären stromerzeugten VOA zu einem stark lokalen konformen Netz ist durch entsprechende Voraussetzungen und Sätze abgesichert \cite{cklw}.

\key{Die zusätzliche Ursprungshypothese ist klar lokalisierbar: Die physische TFPT-Quelle ist diese minimale positive lokale Level-1-Vakuumquantisierung des ursprünglichen markierten Gitters. Innerhalb dieser Hypothese sind Vakuum, Ladungsfelder und konforme Zeit gemeinsam festgelegt.}

Die direkte Gitterroute ist bereits in den Originalen angelegt, insbesondere in \file{v459_seam_lattice_voa_route.py} und der Matter-Geometry-Route. Ein noch fehlender spezieller CAR-/QWZ-Skalierungsnachweis ist deshalb kein universelles Hindernis für jede Konstruktion der Quelle. Er bleibt erforderlich, wenn genau dieser mikroskopische Ursprung behauptet wird.

\subsection{$C_4+X_{20}$: Der kleine Erzeugersatz}
Die ursprünglichen Felder sind
\begin{equation}
C_a,\quad X_{ia},\qquad a=0,1,2,3,\quad i=0,1,2,3,4,
\end{equation}
zusammen mit ihren Adjungierten. Dabei transformiert $X$ als $\overline{\mathbf5}\otimes\mathbf4$, $C$ als $\mathbf1\otimes\mathbf4$. In der ursprünglichen Wurzelbasis ist P2 diagonal; in der nativen Prozessbasis muss derselbe Generator passend transformiert werden. Seine Wirkung ist
\begin{equation}
Q(Y)=-\sum_{ij}Y_{ij}S_{ji},\qquad Q(Y)|_X=-Y^T\otimes I_4,
\qquad Q(Y)|_C=0.
\end{equation}
Die Sitzung hat diese Zuordnung auf allen 24 ursprünglichen Feldern geprüft.

Die vierzig Wurzeln aus $X$ und $X^\dagger$ schließen auf 72 Wurzeln und acht Cartanrichtungen: die $80$-dimensionale $A_8$-Stromalgebra. Mit den acht Wurzeln aus $C$ und $C^\dagger$ wächst der Abschluss dagegen
\begin{equation}
48\longrightarrow140\longrightarrow240\quad\text{Wurzeln},
\qquad240+8=248.
\end{equation}
Es werden somit keine fremden Erzeuger nachträglich angefügt. Die ursprünglichen vier $C$-Felder liefern genau die Vervollständigung, die dem $X$-Block allein fehlt.

\begin{center}
\begin{tikzpicture}[node distance=9mm]
\node[box,text width=4cm] (x) {$X_{ia}$ und $X_{ia}^\dagger$\\20 Felder plus Adjungierte};
\node[box,right=of x,text width=3cm] (a8) {$A_8$\\72 Wurzeln + 8 Cartan};
\node[box,right=of a8,text width=3cm] (e8) {$E_8$\\240 Wurzeln + 8 Cartan};
\node[box,below=of a8,text width=3cm] (c) {$C_a,C_a^\dagger$\\ursprünglicher Viererblock};
\draw[arr] (x)--(a8);\draw[arr] (a8)--(e8);\draw[arr] (c)--(e8);
\end{tikzpicture}
\end{center}

\picturetext{Die zwanzig Felder erzeugen mit ihren Adjungierten bereits eine große geschlossene Instrumentengruppe. Vier weitere, schon ursprünglich vorhandene Stimmen verbinden sie mit den fehlenden Klangrichtungen. Danach ist das gesamte Instrument erreichbar.}

\subsection{Die $A_8$-Erweiterung und ihre Grenzen}
Mit ursprünglichen $X$-Wurzeln lässt sich eine einfache $A_8$-Basis konkret schreiben:
\begin{align}
\beta_i&=X_{i+1,0}-X_{i,0} &&(i=0,\ldots,3),\\
\beta_4&=-X_{4,0},\\
\beta_{5+a}&=X_{0,a}-X_{0,a+1} &&(a=0,1,2).
\end{align}
Hier bezeichnet $X$ in der Gleichung den Ladungsvektor des entsprechenden Feldes. Die Gramdeterminante ist $9$; die Smith-Normalform der Inklusion besitzt sieben Einträge $1$ und einen Eintrag $3$. Der \emph{Gitterquotient} ist deshalb $E_8/A_8\cong\Z_3$. Die Wurzeln zerfallen in $72+84+84$ und die Ströme in
\begin{equation}
\mathbf{248}=\mathbf{80}\oplus\mathbf{84}\oplus\overline{\mathbf{84}}.
\end{equation}
Für die verwendete Orientierung liegen die $C$-Wurzeln im $\Lambda^6$-Sektor. Die $SU(9)/\Z_3$-Stromalgebraerweiterung besitzt bekannte Vorläufer \cite{distler}; neu in dieser Sitzung ist ihr expliziter Anschluss an die ursprünglichen markierten TFPT-Felder und die ausführbare Antwortkette.

Der Gitterquotient $\Z_3$ ist nicht ohne Weiteres die Familienclock $\sigma_F$. Gleiche Gruppenordnung ersetzt keine identische Wirkung. Die noch zu prüfende Gleichsetzung muss auf denselben Feldern, Phasen und Zuständen formuliert werden.

\section{Code, Quartik und Rekursion: Was die endlichen Modelle wirklich zeigen}\label{sec:finite}

\subsection{Vom Binärcode zu den 60 Ereignissen}
Der verwendete binäre Code ist
\begin{equation}
C=RM(1,3)=[8,4,4],\qquad W_C(y)=1+14y^4+y^8.
\end{equation}
Unter der Gaussian-$\mu_4$-Wirkung werden die 240 reellen $E_8$-Wurzeln zu sechzig komplexen Strahlen. Ihre Projektoren lassen sich als gemeinsame Pauli-Eigenprojektoren darstellen. Die fünfzehn primitiven Kontexte enthalten je vier Strahlen. Andere orthogonale Gruppierungen erzeugen dadurch nicht automatisch neue primitive Kontexte.

Zu jedem Strahl gehört die Reflexion
\begin{equation}
r_\ell=I_4-2\ket{\psi_\ell}\!\bra{\psi_\ell},
\qquad \bar U_4=\frac1{60}\sum_\ell r_\ell^{\otimes4},
\qquad G_4=\frac35I-\bar U_4.
\end{equation}
Dabei ist $\bra{\psi}$ die adjungierte Zeile von $\ket\psi$. Das Spektrum von $\bar U_4$ lautet, mit Multiplizitäten als Hochzahlen,
\begin{equation}
(-1)^1,\quad(-1/5)^{45},\quad(1/15)^{135},\quad(1/5)^{70},\quad(3/5)^5.
\end{equation}
Der obere Fünferraum ist durch eine Lücke $2/5$ getrennt. Auf diesem Code gilt ein echter Intertwiner
\begin{equation}
r_\ell^{\otimes4}V=VT_{ab}.
\end{equation}
Sechzig Ereignisse werden auf fünfzehn logische Transpositionen abgebildet, mit vier Ereignissen pro Faser. Die Transpositionen, perfekten Matchings, Pauli-Kontexte und Hecke-Marken haben zwar jeweils fünfzehn Elemente, sind jedoch unterschiedliche Objekte. Nur die tatsächlich konstruierten Abbildungen verbinden sie.

\subsection{Warum die Quartik einen Fünfer erzeugt}
Für die konkrete Pauli-Gruppe der Ordnung 64 wurde der Molien-Ausdruck
\begin{equation}
M(t)=\frac{1-t^{16}}{(1-t^4)^5}
\end{equation}
ausgewertet. Die Dimensionen in den Graden $0,4,8,12,16$ sind $1,5,15,35,69$. Somit ist der Fünfer der erste nichtkonstante Pauli-invariante Polynomraum, \emph{wenn} der kleinste positive Grad als Auswahlregel verlangt wird. Die Auswahlregel selbst ist eine zusätzliche Annahme, solange ihre Herkunft nicht gezeigt ist.

Die projizierte Vierfachpräparation $V^\dagger x^{\otimes4}$ liefert sechs Nullsummenkoordinaten auf der Igusa-Quartik:
\begin{equation}
\sum_i z_i=0,\qquad
\left(\sum_i z_i^2\right)^2=4\sum_i z_i^4.
\end{equation}
Dies ist eine spezielle Präparationsfläche im logischen Fünfer. Ein beliebiger Zustand dieses Fünfers liegt nicht notwendig auf ihr. Die fünf Quartikpolynome der Abbildung sind Pauli-invariant und bilden zusammen einen $G_{31}$-Kovarianten. Die Igusa-Quartik ist deren quartische Relation in den Zielkoordinaten. Beides ist von den skalaren $G_{31}$-Invariantengraden $8,12,20,24$ zu unterscheiden.

Die Projektoren erfüllen $P=SQ_P$ und
\begin{equation}
\rank(S,Q_P,P,S-P)=(35,16,5,30).
\end{equation}
Paarablesungen erfassen zunächst zehn Operatorrichtungen. Mit den ausdrücklich zugelassenen aktiven Paarkontrollen wächst der Beobachtungsraum $10\to20\to25$. Vollständige Beobachtbarkeit ist hier ein Ergebnis der gesamten Kontrollfamilie; sie folgt nicht schon aus einer passiven Paarablesung.

\subsection{Der Vierertensor und die 60/40-Zerlegung}
Der symmetrische Delta-Tensor $A$ und der ursprüngliche Simplextensor $C$ liefern
\begin{equation}
E=A-2C,\qquad F_I=4C-A/3,\qquad G=C+A/6.
\end{equation}
Die normierten Richtungen $\widehat E$ und $\widehat F_I$ sind orthogonal, und
\begin{equation}
\widehat G=\sqrt{3/5}\,\widehat E+\sqrt{2/5}\,\widehat F_I.
\end{equation}
Die 60/40-Aussage bezeichnet also die quadrierten Normanteile dieser beiden Richtungen. Sie ist keine universelle Zerlegung physischer Energien. Als $125\times5$-Matrix gelesen liefert $G_{abcd}/\sqrt2$ den ursprünglichen Dreierencoder.

\subsection{Zwei unterschiedliche, jeweils sinnvolle Rekursionen}
Die ältere gemeinsame Ereignisbindung ist
\begin{equation}
k=I-\frac1{15}\sum_{ab}T_{ab}\otimes T_{ab}.
\end{equation}
Sie besitzt das Spektrum $0^1,(2/5)^5,(2/3)^9,(4/5)^{10}$ und den eindeutigen neutralen Zustand $\Omega_5$. Werden die beiden Ereignislabels unabhängig gezogen, entsteht stattdessen die skalare Energie $16/25$. Die gemeinsame Ereignisidentität ist deshalb eine reale Voraussetzung.

Das Dreieck $k_{12}+k_{23}+k_{31}$ besitzt einen Fünfergrundraum bei Energie $4/5$ und eine Lücke $2/5$. Das ist ein exakter rekursiver Codezweig. Das räumliche Inzidenznetz ist jedoch dreieckfrei. Deshalb darf derselbe Dreiecksencoder nicht ohne neue Identifikation zur lokalen Raumregel erklärt werden.

Der spätere ladungskovariante Pfad verwendet dagegen
\begin{equation}
h_{\rm cov}=\frac56(I-P_\Omega),\qquad
(Wx)_{abc}=\frac{\delta_{ab}x_c+x_a\delta_{bc}}{\sqrt{12}}.
\end{equation}
Er wirkt auf $U\otimes\overline U\otimes U$ und erfüllt
\begin{equation}
(X_1-X_2^T+X_3)W=WX.
\end{equation}
Sein Grundraum ist fünfdimensional, seine Grundenergie $2/3$ und seine Lücke $1/3$. Der Ladungsgenerator wird wirklich durch die Rekursion transportiert.

Vollständige kollektive native $S_6$-Symmetrie zusammen mit Erhaltung des tatsächlichen P2-Markers verbindet die $5$-, $9$- und $10$-Sektoren und erzwingt die Form
\begin{equation}
H=cI+\kappa(I-P_\Omega).
\end{equation}
Positivität und ein eindeutiger neutraler Nullzustand verlangen $\kappa>0$ nach entsprechender Offsetwahl. Der konkrete Wert $5/6$ benötigt zusätzlich die erhaltene ursprüngliche Spur. Eine mögliche Instrumentenoperation ist dabei nicht automatisch eine Symmetrie des Hamiltonoperators.

\subsection{Was bei der Blockkomposition erhalten bleibt}
Die positive Neun-Port-Vervollständigung erfüllt in beiden Orientierungen
\begin{equation}
H_9V=V(4I+h).
\end{equation}
Auf einem vorgegebenen endlichen Blockgraphen bleibt das Produkt der Codes daher unter einer Summe solcher Links invariant: $PHQ=0$. Die ursprünglichen einzelnen schwachen Restlinks haben dagegen Leakage. Die Neun-Port-Regel ist somit eine zusätzliche, ausdrücklich benannte Kompositionsregel. Wenn die Graphstruktur selbst dynamisch variiert, darf der Offset $4I$ nicht pauschal als bedeutungslos entfernt werden.

Der ältere Quartikencoder zerlegt sich gegenüber dem Pfadencoder als
\begin{equation}
V_G=\sqrt{2/3}\,W+\sqrt{1/3}\,R,
\qquad P_G(t)=\frac{5+4\cos t}{9}.
\end{equation}
Unter beliebigen kollektiven Ladungsoperationen wächst der gewählte Restfünfer auf den gesamten 70er-Sektor. $5+70=75$ ist deshalb der nächste notwendige Träger dieser Operationen. Weitere ursprüngliche Operationen erreichen auch den 45er-Sektor. Die vollständige Quelle muss diese Fortsetzungen bewahren; ein kleiner Restträger kann nicht allein wegen seiner Bequemlichkeit geschlossen erklärt werden.

\section{Die gemeinsame Zeit: Der alte und der neue Quellblock passen wirklich zusammen}

Die beiden Beschreibungen $D_5+A_3$ und $A_4+A_3+\mathfrak u(1)$ sind jetzt am gleichen achtfachen Heisenbergraum verbunden. Mit $J_5$ als Fünfermatrix aus lauter Einsen lauten die Projektoren
\begin{align}
P_{A_4}&=\diag(I_5-J_5/5,0_3),\\
P_{A_3}&=\diag(0_5,I_3),\\
P_{\mathfrak u(1)}&=\diag(J_5/5,0_3).
\end{align}
Sie summieren sich zu $I_8$, und $P_{A_4}+P_{\mathfrak u(1)}=P_{D_5}$. Deshalb ist der tatsächliche Virasorovektor
\begin{equation}
\omega=\frac12\sum_{j=1}^8h_j(-1)^2\Omega
\end{equation}
in beiden Zerlegungen gleich. Damit gilt auf dieser Quelle
\begin{equation}
T_{E_8}=T_{A_8}=T_{D_5}+T_{A_3}
=T_{A_4}+T_{A_3}+T_{\mathfrak u(1)}.
\end{equation}
Die Zentralzahlen erfüllen $8=5+3=4+3+1$, aber diese Zahlen sind nur eine Konsequenz des stärkeren Projektorbeweises. Gleich ist der Stress-Tensor samt Moden, nicht lediglich seine Zentralzahl.

\begin{center}
\begin{tikzpicture}
\node[box,text width=4.8cm] (left) {$D_5$\quad +\quad $A_3$\\$5$\quad +\quad $3$};
\node[box,right=18mm of left,text width=6.2cm] (right) {$A_4$\quad +\quad $\mathfrak u(1)$\quad +\quad $A_3$\\$4$\quad +\quad $1$\quad +\quad $3$};
\draw[arr] (left)--node[above,font=\small]{derselbe $T$}(right);
\node[below=7mm of right,align=center,font=\small] {Die acht Richtungen werden anders gruppiert.\\Die gemeinsame konforme Zeit bleibt gleich.};
\end{tikzpicture}
\end{center}

Auf den tatsächlichen ursprünglichen Feldern ergibt sich
\begin{equation}
h(X)=\frac25+\frac38+\frac9{40}=1,
\qquad h(C)=0+\frac38+\frac58=1.
\end{equation}
Die ergänzende normierte $\mathfrak u(1)$-Richtung ist
$(1/\sqrt5,\ldots,1/\sqrt5,0,0,0)$; ihre Ladungen sind $3/(2\sqrt5)$ auf $X$ und $-\sqrt5/2$ auf $C$. Sie steht orthogonal zur P2-Hyperladung. Ihre Benennung als $U(1)$ erlaubt daher keine Gleichsetzung mit Elektromagnetismus oder Hyperladung.

Die geometrische modulare beziehungsweise konforme Zeit ist innerhalb des Vakuumnetzes bestimmt. Ein interner Clock-Zeiger, eine diskrete Survivalmatrix, $L_0$, eine räumliche Translation und die Laborzeit sind trotzdem unterschiedliche Objekte. Die physische Realisierung muss ihre passende Verflechtung liefern.

\section{Das ausführbare Quellgesetz und die echte Hyperkante}\label{sec:ward}

\subsection{Eine Regel für alle endlichen Stromwörter}
Für $F_n=\omega(J_{A_1}(z_1)\cdots J_{A_n}(z_n))$ bestimmt die affine Ward-Rekursion
\begin{equation}\label{eq:ward}
F_n=\sum_{j=2}^n\frac{\kappa(A_1,A_j)}{(z_1-z_j)^2}
F_{n-2}^{\widehat1,\widehat j}
+\sum_{j=2}^n\frac{1}{z_1-z_j}F_{n-1}^{[1,j]},
\quad F_0=1,\quad F_1=0.
\end{equation}
Im zweiten Term wird $A_j$ durch $[A_1,A_j]$ ersetzt. Alle Zwischenströme bleiben in der wirklichen 248-dimensionalen Wurzel-/Cartanalgebra. Die Implementierung nutzt exakte rationale Positionen und rationale Arithmetik; sie ist keine Näherung durch Gaußpaarungen.

Die öffentlichen Funktionen \file{evaluate_source_word} und \file{connected_source_word} nehmen $C$-/ $X$-Felder, ihre Adjungierung und paarweise verschiedene endliche Positionen entgegen. Der Kern ist für beliebige endliche Wörter formuliert; die Webschnittstelle begrenzt die Wortlänge aus Laufzeitgründen auf acht. Diese Grenze ist keine Aussage über die Physik. Exakte Ergebnisse werden als Bruch erhalten; ein nicht darstellbarer Gleitkommawert wird nicht als gültige Zahl ausgegeben.

\subsection{Eine vollständig verbundene Sechserantwort}
Für die Folge
\begin{equation}
X_{00},\ X_{10}^{\dagger},\ X_{11},\ X_{21}^{\dagger},\ X_{22},\ X_{02}^{\dagger}
\end{equation}
ist die Gesamtladung null. Jede der $2^6-2=62$ echten nichtleeren Teilmengen besitzt dagegen eine nichtverschwindende Gesamtladung. Alle ihre Vakuumkorrelatoren verschwinden, ebenso alle fünfzehn vollständigen Zweierpaarungsprodukte. Die ganze Antwort ist trotzdem
\begin{equation}\label{eq:hex}
F_6=\frac{1}{(z_1-z_2)(z_1-z_6)(z_2-z_3)
(z_3-z_4)(z_4-z_5)(z_5-z_6)}.
\end{equation}
Bei $(-5,-2,-1,1,3,7)$ ist der Wert exakt $1/576$. Er wurde unter allen 720 gleichzeitigen Permutationen von Feldern und zugehörigen Positionen verglichen. Wird nur die erste Position auf $-4$ verschoben, ergibt sich $1/352$; bei $(-4,-1,0,2,5,9)$ ergibt sich $1/936$.

\begin{center}
\begin{tikzpicture}[scale=1.0]
\fill[teal!9] (90:1.9)--(30:1.9)--(-30:1.9)--(-90:1.9)--(-150:1.9)--(150:1.9)--cycle;
\foreach \ang/\lab in {90/1,30/2,-30/3,-90/4,-150/5,150/6}{\node[circle,draw=teal,fill=white,minimum size=7mm] (v\lab) at (\ang:1.9) {\lab};}
\node[align=center,font=\small] at (0,0) {ganze Antwort\\$1/576$};
\node[align=left,text width=7.2cm,right=20mm of v2] {\textbf{Alle sechs gemeinsam zählen.}\\Jede echte Teilgruppe ist ladungsungleichgewichtig und hat Antwort null. Die ganze Gruppe ist neutral und antwortet.};
\end{tikzpicture}
\end{center}

Mit
\begin{equation}
Z(t)=\omega\!\left(\prod_i(1+t_iJ_i(z_i))\right)
\end{equation}
ist der verbundene Beitrag der Koeffizient von $t_1\cdots t_n$ in $\log Z$. Weil hier alle echten Teilantworten verschwinden, ist die volle Sechserantwort bereits vollständig verbunden. Ein unabhängiger Vergleich über sämtliche 203 Mengenpartitionen eines anderen gemischten Sechserworts ergab volle Antwort $1/720$ und verbundenen Anteil $-1/8640$. Eine Cartan-Viererantwort liefert den passenden Gauß-Gegenfall: volle Antwort ungleich null, verbundener Anteil null.

\key{Diese Hyperkante ist ein tatsächlich berechneter Koeffizient derselben Quelle. Sie ist kein zusätzlich postuliertes Sechskörper-Hamiltontermchen und beweist keine nichtlokale Signalübertragung.}

\subsection{Die $W$-Rekursion ist eine besondere Antwort dieses Gesetzes}
An den harmonischen Marken $(-1,0,1,\infty)$ ergibt der vollständige $X$-Vierertensor
\begin{equation}
F_{ia,jb,kc,ld}=(\delta_{ij}\delta_{kl}+\delta_{il}\delta_{jk})
(\delta_{ab}\delta_{cd}+\delta_{ad}\delta_{bc}).
\end{equation}
Definiert man
\begin{equation}
(N_d)_{abc,x}=\delta_{ab}\delta_{cx}+\delta_{ax}\delta_{bc},
\qquad W_d=\frac{N_d}{\sqrt{2(d+1)}},
\end{equation}
so ist $W_d^\dagger W_d=I_d$. Nach der passenden Anordnung der Tensorfaktoren folgt
\begin{equation}
W_{\rm gemeinsam}=W_5\otimes W_4,
\qquad N_5^\dagger N_5=12I_5,\quad N_4^\dagger N_4=10I_4,
\quad N_{\rm gemeinsam}^\dagger N_{\rm gemeinsam}=120I_{20}.
\end{equation}
Der mittlere Faktor ist jeweils konjugiert. Alle 8000 Ausgabetripel der dünn besetzten Abbildung wurden geprüft. Bei festem Familienindex wird derselbe Vierertensor zu $2(\delta_{ij}\delta_{kl}+\delta_{il}\delta_{jk})$. Mit der tatsächlichen reflexionspositiven Paarung ist er $K=2I+10P_\Omega$, sodass $I-K/12=h_{\rm cov}$ folgt.

Damit verbindet sich die früher separat motivierte Bindung direkt mit der tatsächlichen Quellenantwort. Der Geltungsbereich bleibt jedoch konkret: Ein gemischtes Wort $C_0,X_{00}^\dagger,X_{01},C_1^\dagger$ liefert $1/2$ statt des naiven Wertes $1$ einer Erweiterung zu $W_6\otimes W_4$. Auch ein naiver einziger $W_{20}$ ersetzt das Produkt nicht.

Im Koaleszenzgrenzwert eines Sechserworts schließt die Antwort auf $(W\otimes\overline W)\Omega$. Bei endlichem Verhältnis $t=\epsilon/L$ ist die Leakage positiv und beginnt mit $t^2/3$. Die Sektoren $5_W+5_Z+70+45$ bleiben deshalb erforderlich. Die affine Gramrechnung zeigt beim Niveau $J_{-1}X$ die Eigenwerte $6^5,2^{45},0^{70}$; beim Niveau $J_{-2}X$ sind diese um eins erhöht. Die Quelle besitzt also einen Nachfahrenturm. Der kleine $W$-Code beschreibt einen kontrollierten Grenzsektor dieses Turms.

\subsection{KZ, Fusionsbäume und echte Komposition}
Die betrachtete $SU(5)$-Level-1-KZ-Verbindung ist flach, obwohl überlappende Residuen nicht miteinander kommutieren. Zwei klassische Tensorkoordinaten bedeuten dabei nicht zwei unabhängige physische Level-1-Fusionsblöcke. Die Partnersektoren des vollständigen $E_8$-Feldes kompensieren die gebrochenen Phasen des isolierten Faktors.

Die rationalen Hecke-Charakterphasen erhalten Wurzelklammer, Paarung, Ward-Regel und diagonalen Casimir. Damit ist ihr Anschluss an die vollständige Quelle real. Eine Hecke-Inklusion wird dadurch nicht zu einem KZ-Transport. Ebenso sind die $E_8$-Ströme Vakuumnachfahren und keine zusätzlich eingeführten adjungierten $E_8$-Level-1-Primärfelder. Analytische Aussagen zu Sewing und Faktorisierung benötigen die entsprechenden Konvergenzvoraussetzungen \cite{gui,guizhang}.

\section{Die Quellen- und Clockprüfungen: Welche Alternativen wirklich ausgeschlossen sind}

\subsection{Die frühe Kanalfamilie und die vollständige Normierung}
Eine frühe Phase der Sitzung betrachtete die endliche Familie $\Phi_b$, $0\le b\le2/9$. Ihr Populationsschatten bleibt gleich, während eine Kohärenz den Eigenwert
\begin{equation}
\lambda_{12}=2b-\frac19
\end{equation}
besitzt. Damit ist die reine Populationsauslese in dieser Klasse unzureichend, um den Kanal zu bestimmen. Die Familie ist kein Gegenmodell zu sämtlichen TFPT-Prinzipien, weil nicht alle ursprünglichen Markierungen und höheren Antworten in ihre Definition eingingen.

Der kohärente Lift erfüllte auf dem Code die Normierung, aber nicht auf dem ganzen 256-dimensionalen Raum:
\begin{equation}
q=\sum_\mu L_\mu^\dagger L_\mu,\qquad
q|_{\rm code}=I,\qquad q|_{\Lambda^4\C^4}=\frac{10}{9}I.
\end{equation}
Die funktionale Reparatur $B_\mu=L_\mu q^{-1/2}$ normiert den gesamten Lift und erhält die Code-Krauszweige. Sie wählt den Parameter $b$ nicht. Dass $q^{-1/2}$ zur endlichen erzeugten Algebra gehört, beweist auch keine einfache lokale physische Realisierung. Der frühere HTML-Gesamtbericht dokumentiert diesen Zwischenstand und darf nicht als letzte Aussage der Sitzung gelesen werden.

\subsection{Die ursprüngliche Clock ist ein konkreter markierter Operator}
Der v487-Survivaloperator lautet
\begin{align}
D&=\begin{pmatrix}1&0&0\\0&0&1\\0&1&0\end{pmatrix},\qquad P_0=\diag(1,0,0),\\
M&=\frac12I+\frac16D+\frac13P_0,\qquad \operatorname{spec}M=(1,2/3,1/3).
\end{align}
Hier ist $D$ der Decktausch von Ordnung zwei, während $\sigma_F$ die Familienmarkierung von Ordnung drei bezeichnet. Seine Eindeutigkeit gilt in der ursprünglichen treuen markierten Dreierklasse. Die Fortsetzung
$C=\frac12I+\frac16S_{13}+\frac13WW^\dagger$ auf den Trimer setzt zusätzliche Identifikationen voraus und bestimmt außerhalb des Seedsektors nicht automatisch alles.

Auch die beiden früher gleichgesetzten Fünferclocks sind verschieden. Der Riemann--Roch-Charakter ist $(5,1,1,1)$, der zyklische Codecharakter $(5,3,5,3)$. Der Hom-Raum ist neundimensional, sein maximaler Intertwinerrang aber nur drei. Es gibt also Intertwiner, jedoch keinen invertiblen markierten Intertwiner zwischen diesen beiden Fünferwirkungen. Ein gemeinsamer Siebenerträger repariert nur diese isolierten $C_4$-Wirkungen.

Der volle unmarkierte Cartan-$M_4$-Anschluss scheiterte an der ursprünglichen $\sigma_F$-Markierung: Die Fixalgebren $M_2\oplus\C\oplus\C$ und $M_2\oplus M_2$ besitzen Dimension sechs beziehungsweise acht. Die korrekte native Fixebene ist $\C^2$, mit zwölf Wurzeln und drei $\mu_4$-Strahlen. Dort ergibt $E_i=(2/3)P_i$ die ursprüngliche Trine.

Ihre Erwartungswertmatrix
\begin{equation}
B=\frac1{18}\begin{pmatrix}13&1&4\\1&13&4\\4&4&10\end{pmatrix}
\end{equation}
ist keine Übergangsmatrix aufeinanderfolgender gleicher Trinemessungen: Jede einzelne Ausgangswahrscheinlichkeit einer solchen Messung ist höchstens $2/3$, während $B_{11}=13/18$ größer ist. Dies ist ein kleiner, entscheidender Gegenfall gegen die falsche Interpretation.

\subsection{Die korrekt markierte Ereignisfaser}
Native Reflexionswörter realisieren die registrierte v976-Faser
\begin{equation}
w_t=(1/2+t,t,t,1/18-t,2/9-t,2/9-t),
\quad q=6t,\quad0\le t\le1/18.
\end{equation}
Die GNS-Antwort mit linken Einsetzungen ist $FAAF=q/9$. Die ältere Vorgabe $1/27$ wählt deshalb $q=1/3$, sofern ihre unabhängige Quellenidentifikation gilt. Dies ist eine konkrete Möglichkeit, wie eine zusätzliche ursprüngliche Antwort eine zuvor offene Faser auswählt.

Für jede Faser ist der stationäre reversible $S_3$-Random-Walk mit Gleichverteilung $1/6$ und Faltungstransfer $T_q$ bei allen endlichen Historientests reflexionspositiv. Diese Historienfunktion ist nicht automatisch die zuvor genannte linksinsertierte $M_2$-GNS-Antwort. Eine reversible kontinuierliche Einbettung mit nichtnegativen $S_3$-Ereignisraten existiert dagegen genau im geprüften Bereich $2/9\le q\le1/4$. Bei $q=1/3$ ist eine der Raten $\log(3/4)/6<0$. Ausgeschlossen ist somit diese spezielle klassische kontinuierliche Ereignisratenklasse. Eine allgemeine Quanten-Markov-Einbettung ist damit nicht widerlegt.

Die frühere $b$-Familie auf $M_4$ und diese spätere markierte $q$-Faser auf $\C^2$ sind nicht dieselbe Parametrisierung derselben vollständigen Quelle. Ihr jeweiliger Prüfbereich muss erhalten bleiben.

\subsection{Höhere Antworten unterscheiden sogar gleiche Wurzelzerfallsraten}
Die Phasenkanalprüfung betrachtete eine ausdrücklich beschränkte Klasse zufälliger unitärer Generatoren
\begin{equation}
\mathcal L_O(A)=\frac1{|O|}\sum_{m\in O}(U_m^\dagger A U_m-A).
\end{equation}
Die nichttrivialen Charakterorbits haben Größen $15,120,120$. Ein Deckorbit und ein Braidorbit liefern bei derselben gesamten Sprungrate auf allen 240 Wurzelströmen dieselbe Zerfallsrate $16/15$. Auf einem Grad-zwei-Vertex mit $\beta=(1+i)\alpha$ und $\norm\beta^2=4$ unterscheiden sie sich dagegen: Rate null beziehungsweise $16/15$.

Das ist ein exakter Unterschied in einer höheren Zukunftsantwort. Er bewahrt unter anderem die untersuchte native Symmetrie, Ladungsaddition und Vakuumkovarianz. Er erzwingt aber nicht schon die ursprüngliche primitive Ereignisauswahl, alle Clockantworten, sämtliche registrierten Mehrzeitwörter oder die physische Raumzeitabbildung. Die korrekte Aussage lautet deshalb: \emph{Nicht-Eindeutigkeit innerhalb dieser Phasenkanalklasse}. Daraus folgt weder eine Mehrdeutigkeit des festgelegten Gittervakuums noch eine Nicht-Eindeutigkeit nach sämtlichen TFPT-Constraints.

\section{Hecke, Double Cover und geladener Carry: Was bei einer Vergröberung mitwandern muss}\label{sec:carry}

\subsection{Die fünfzehn Marken sind an die ursprünglichen Ereignisse angeschlossen}
Mit $\pi=1+i$ gilt im verwendeten Gaussian-Gitter
\begin{equation}
L/\pi L\cong\mathbb F_2^4,\qquad
M_y=\{x:\overline h(x,y)=0\},\qquad
\chi_y(x)=(-1)^{\overline h(x,y)}.
\end{equation}
Alle sechzig nativen Ereignisse wurden tatsächlich auf die fünfzehn Marken abgebildet, einschließlich der Äquivarianzprüfung. Vier Ereignisse pro Marke bleiben dennoch verschiedene Operationen.

Der konkrete Deckcharakter mit Kennung 152 liefert ein $D_8$-Untergitter von Index zwei: 112 gerade Wurzeln und 128 ungerade Spinorwurzeln. $V_{D_8}$ zusammen mit seinem tatsächlichen Spinorkoset stellt dieselbe $E_8$-Quelle wieder her. Es wird kein unabhängiges neues Materiefeld erfunden.

Die Projektion auf den geraden Teil ist jedoch kein Liehomomorphismus. Für eine ungerade Wurzel gilt
\begin{equation}
P[e_\alpha,e_{-\alpha}]=[e_\alpha,e_{-\alpha}]
\ne0=[Pe_\alpha,Pe_{-\alpha}].
\end{equation}
Die zunächst entfernten Felder können gemeinsam eine erhaltene Cartanantwort erzeugen. Das ist genau der Mechanismus, wegen dessen eine heute unsichtbare Information für spätere Prozesse erhalten bleiben muss.

Beim ersten echten Untergitter sinkt der Rang der geerbten reduzierten Paarung von vier auf zwei. Das kanonische Dual stellt weitere Tests bereit; es erzeugt nicht automatisch zusätzliche Vertexfelder. Über die ersten Kinder ergeben sich die vorhandenen 256 Vorzeichencharaktere.

